% =============================================================
% main.tex — Tiyatro Sistemi Test Dosyası
%
% Derleme: XeLaTeX ile derleyin (pdfLaTeX çalışmaz)
%   xelatex main.tex
%
% TeXstudio: Options → Configure → Build → Default Compiler → XeLaTeX
% =============================================================

\documentclass[draft]{tiyatro}
% Final mod için: \documentclass[final]{tiyatro}

% ── Kapak bilgisi ─────────────────────────────────────────────
\title{Test Oyunu}
\author{Salih Bektaş}
\draft{0.1.0}
\copyrightnotice{2026 Nisan Test}
\address{İstanbul}
\email{test@test.com}

% ── Persona tanımlamaları ─────────────────────────────────────
% \karakter[Tam Ad][Açıklama]{Script Adı}{alias}
% \karakter* → starred: kisiler tablosunda öne çıkar, kapakta gösterilir

\karakter[Maria][Leningrad Komsomol'ü Sekreteri]{MARİA}{ma}
\karakter[Mustafa][\textit{Buralı}, Eski-Devrimci]{MUSTAFA}{mu}
\karakter[Morris][Şikago'lu bir Sovyet gazetecisi]{MORRİS}{must}

% Grup karakter örneği
\karaktergrubu[Pravda  Gazetesi]{
	\gkarakter[Anlatıcı]{ANLATICI}{anl}
	\gkarakter[Kadın Anlatıcı]{KADIN ANLATICI}{kanl}
	\gkarakter[Mikail]{MİKAİL}{mi}
}



% ── Custom font (opsiyonel) ───────────────────────────────────
% Proje klasöründe OzelFont.ttf varsa:
%   \fontyukle{taslakfont}{OzelFont.ttf}
% Yoksa bu satırı yorum yapın, kapakta varsayılan font kullanılır

\begin{document}

\maketitle      % Kapak sayfası
\kisiler        % Kişiler tablosu (dramatis personae)
\pagebreak

% =============================================================
% PERDE I
% =============================================================

\perde[Birinci Perde]

% ── Sahne 1 ───────────────────────────────────────────────────
% Yer/zaman bilgisiyle
\sahne[İç, Apartman — Gece]

% Sahne yönü — isimsiz, karakter bağımsız
\action{Işıklar yavaşça açılır. Sahne boştur.}

% Diyalog
\mu Merhaba! Buralara hep gelir misiniz?

% Diyalog içinde inline eylem — \ic{} komutu
\ma Selam. \ic{gülümser} Nasılsın?

\mu İyiyim, teşekkürler.

% Diyalog içinde karakter referansı — Persona özelliği
% \ma burada sadece isim basar, yeni diyalog başlatmaz
% Ayşe'nin adı \karakter tanımından değiştirilirse burada da değişir
\mu \ma nerede diye soruyordum aslında.

\ma Ben buradayım zaten!

% Altı çizili vurgu — _kelime_ sözdizimi
\mu Ama sen _orada_ duruyordun!

% ── Sahne 2 ───────────────────────────────────────────────────
% Manuel numara örneği — sayaç kendi yolunda devam eder
\sahne[Dış, Sokak — Sabah][1.1]

% Zincirleme action — birden fazla karakter aynı anda
% Çıktı: (Ali, Ayşe, Mustafa: Girerler.)
\mudo\mado\mustdo Girerler.

\must Günaydın efendim.

\mu Günaydın \must.

\ma \ic{şaşırır} Sen de mi geldin?

\mustdo Başını eğer.

% Çok dilli kullanım
\mu \ru{Добрый день} dedi ve güldü.

\ma \en{Good morning!} diye yanıtladı.

% =============================================================
% PERDE II
% =============================================================

\perde[İkinci Perde]

\sahne[]

\mu İkinci perdedeyiz artık.

\ma Zaman çabuk geçti.


\end{document}
